\chapter{关键算法实现与优化}

\section{利用环结构降低采样开销}
\subsection{大矩阵从哪里产生}
Sageloli的发钥和签名都需要在格陪集中采样。相关扰动由陷门决定其形状，再与离散采样配合，得到满足方程和范数要求的短向量。实现既要得到正确的相关性，也要承受论文参数下的对象规模。

环元素是一组多项式系数，乘法可以直接在负循环环上计算。如果把每个环元素展开成旋转矩阵，后续陷门与协方差也会随之变成很大的实数矩阵。PARAM-I取$d=2048$、$\ell=132$、$k=65$，签名相关扰动的系数维数为
\[
 D=d(\ell+k)=403456.
\]
仅显式保存一个$D\times D$实数矩阵，按每项8字节估算就需要约1.18 TiB，还未计入分解工作区。因此，实际实现需要利用原式的结构来组织计算。

本作品始终保存多项式系数。公钥、主陷门和身份陷门分别以$\ell\times d$、$(\ell-k)\times k\times d$和$\ell\times k\times d$的系数数组表示。环乘法在原环上执行，结果仍检查
\[
 A=[\bar A\mid g-\bar AR],\qquad AT_s=H_1(P,\rho,s)-g\pmod q.
\]
这样，公钥和密钥关系保持不变，计算过程避免展开旋转矩阵。

\subsection{把大分解转为逐频率的小分解}
设当前陷门为$T$，其形状为$\ell\times k$个环元素，令
\[
 J=\begin{bmatrix}-T\\I_k\end{bmatrix},\qquad
 a=s^2-r^2,\qquad c=9r^2.
\]
这里$r$是基础采样宽度，$s$是该次相关扰动使用的宽度。需要计算的协方差形状为
\[
 \Sigma=aI-cJJ^\star,
\]
其中$\star$表示环上的伴随。这个式子的主要部分是标量矩阵$aI$，修正部分由只有$k$列的$J$产生。利用这一点，可以将分解限制在较小的子空间。

第一步，对多项式执行负循环Fourier变换。在每个频率$\nu$上，环乘法变为复数乘法，$J$对应一个$(\ell+k)\times k$的复数矩阵$J_\nu$。第二步，计算Gram矩阵
\[
 G_\nu=J_\nu^*J_\nu
      =U_\nu\operatorname{diag}(\lambda_i)U_\nu^*.
\]
$G_\nu$只有$k\times k$；在PARAM-I中，最大特征分解为$65\times65$。

\begin{figure}[H]
\centering
\begin{tikzpicture}[node distance=8mm and 7mm]
 \node[flow,text width=39mm] (poly) {陷门多项式$T$\\保存系数数组};
 \node[flow,text width=39mm,right=of poly] (fourier) {负循环Fourier变换\\$J_\nu$：$197\times65$};
 \node[flow,text width=39mm,right=of fourier] (gram) {计算$J_\nu^*J_\nu$\\分解$65\times65$矩阵};
 \node[check,text width=39mm,below=of gram] (apply) {对新随机向量应用\\低秩平方根};
 \node[check,text width=39mm,left=of apply] (inverse) {逆变换回系数域\\逐系数离散采样};
 \node[check,text width=39mm,left=of inverse] (check) {检查精确模方程\\和原方案范数条件};
 \draw[flowarrow](poly)--(fourier);
 \draw[flowarrow](fourier)--(gram);
 \draw[flowarrow](gram)--(apply);
 \draw[flowarrow](apply)--(inverse);
 \draw[flowarrow](inverse)--(check);
\end{tikzpicture}
\caption{结构化采样的计算步骤；矩阵形状对应PARAM-I}
\label{fig:ring-sampler}
\end{figure}

第三步，利用小矩阵的特征分解构造平方根算子
\[
 L_\nu=\sqrt a I-
 J_\nu U_\nu\operatorname{diag}\!\left(
 \frac{c}{\sqrt a+\sqrt{a-c\lambda_i}}
 \right)U_\nu^*J_\nu^*.
\]
该式满足$L_\nu L_\nu^*=aI-cJ_\nu J_\nu^*$。对新随机向量应用$L_\nu$时，按右侧的小矩阵与向量乘法逐步计算，无需形成完整的$L_\nu$。随后逆变换回系数域，继续执行原方案要求的离散Gaussian和三进制陪集采样。方法使用格陷门采样中的扰动思想~\cite{micciancio2012trapdoors}，并针对本方案的环结构组织执行。

需要保存的主要数据量为$O(d\ell k+dk^2)$，对应陷门系数和各频率的小矩阵计划。相较$O(D^2)$的稠密表示，这直接降低了表示和分解规模。已经计算好的计划可在同一陷门下复用，每次发钥和签名仍生成新的随机样本。

\subsection{数值计算与精确方程怎样衔接}
Fourier步骤使用数值运算，方案中的模方程则必须精确成立。实现将两部分分开：数值计划在建立时检查正定充分条件和变换误差；回到系数域后，环方程由整数运算检查。大模数$q=3^k$的公开乘积使用FLINT整数多项式，私有小整数乘法采用双素数NTT与CRT恢复，并检查整数卷积幅度落在可唯一恢复的范围内。

签名生成后继续执行原方案的范数门与验签方程。这样，浮点计算的职责是产生相关扰动，精确算术的职责是验证代数关系；两者在采样接口处衔接。第5章给出PARAM-I与PARAM-II的实际运行成本。

\section{按取值范围编码密钥与签名}
\subsection{JSON表示中的重复开销}
SRGG的公钥、陷门和签名含有大量多项式系数。JSON将每个整数写成十进制字符串，还需要负号、逗号、括号和字段名。对大对象，表示开销会累积到与数学数据相近甚至更大的规模。

紧凑编码依据字段的数学取值范围选择位数。以取值为$-1,0,1$的陷门系数为例，三种可能值只需2 bit；公开模$Q=15744063$的系数落在$[0,Q)$，使用24 bit。对短密钥和签名响应，先以Zigzag映射处理正负号，再选取能够容纳该字段的最小统一位宽。

\begin{table}[H]
\centering\small
\caption{SRGG不同字段的编码方法}
\label{tab:codec-fields}
\begin{tabularx}{\textwidth}{@{}p{.23\textwidth}p{.25\textwidth}X@{}}
\toprule
\rowcolor{black!8} 数据 & 存储方法 & 解码时检查 \\
\midrule
公开模$Q$系数 & 每系数24 bit & $0\leq x<Q$ \\
三值陷门系数 & 每系数2 bit & 拒绝第四个未使用码字 \\
短密钥、响应整数 & Zigzag与最小统一位宽 & 数值范围、位宽最小性 \\
二元挑战 & 每系数1 bit & 固定多项式次数与长度 \\
类型、参数与上下文 & 固定头部和长度前缀 & 对象类别、公钥、参数一致 \\
\bottomrule
\end{tabularx}
\end{table}

对含$N$个系数、统一位宽为$b$的字段，其主要编码长度为$\lceil Nb/8\rceil$字节，再加必要头部。解码后得到与JSON格式相同的数学对象，签名和验证继续使用原算法。

\subsection{四类对象与规范解码}
该方法同时用于公钥、主状态、用户密钥和签名。各类对象有独立类型标识，系数顺序固定，末字节未使用位必须为零。解析器先根据形状计算长度，再读取数据，并检查参数、身份、公钥上下文及取值范围。解码后重新编码应得到相同字节串，从而保证同一个对象只有一种规范表示。

对象尺寸的改善会减少密钥文件和传输字节，实际收益取决于该对象在完整协议中的占比。例如，签名变短可以直接缩短对称加密载荷，但HIBE胶囊和一次性认证仍占固定空间。第5章将四类对象的编码对照与完整密文组成放在一起分析。编解码也有计算成本，实验同时保留该项测量。

\section{通过接收会话减少重复委派}
\subsection{重复计算发生在哪里}
在Lai组合中，接收者持有第一层密钥$\mathsf{dk}_r$。解密来自发送者$s$的消息前，需要运行
\[
 \mathsf{dk}_{r,s}\leftarrow\mathsf{Derive}(\mathsf{dk}_r,s).
\]
同一接收者连续读取同一发送者的消息时，委派输入保持相同；反复执行会重复支付生成叶钥的开销。加密种子、密文与每条消息的签名则各不相同，需要逐条处理。

本作品引入显式接收会话：在会话建立时生成并验证一次叶钥，后续在固定双方身份下复用。图~\ref{fig:receiver-session}展示计算安排的变化。

\begin{figure}[H]
\centering
\begin{tikzpicture}[node distance=6mm and 6mm]
 \node[font=\small,anchor=east] (old) {逐条委派};
 \node[flow,right=of old,text width=30mm] (o1) {委派＋解密1};
 \node[flow,right=of o1,text width=30mm] (o2) {委派＋解密2};
 \node[flow,right=of o2,text width=30mm] (o3) {委派＋解密3};
 \node[font=\small,anchor=east,below=22mm of old] (new) {接收会话};
 \node[check,right=of new,text width=30mm] (prep) {委派一次\\验证并保留叶钥};
 \node[check,right=of prep,text width=30mm] (dec) {解密1、2、3\\逐条认证与验签};
 \node[check,right=of dec,text width=30mm] (close) {关闭会话\\停止使用叶钥};
 \draw[flowarrow](o1)--(o2);\draw[flowarrow](o2)--(o3);
 \draw[flowarrow](prep)--(dec);\draw[flowarrow](dec)--(close);
\end{tikzpicture}
\caption{会话把与密文无关的委派安排到准备阶段}
\label{fig:receiver-session}
\end{figure}

\subsection{身份绑定与会话生命周期}
会话建立时检查公开上下文、接收者身份和一级钥绑定；生成叶钥后检查完整$(r,s)$路径及陷门方程。会话只服务于该公钥和身份对，关闭后停止解密。每条消息继续执行规范解析、外层字节认证、HIBE解封装、GCM验证、IBS验签及一次性公钥承诺核对。

叶钥作为秘密状态保留在进程内，使用者显式关闭会话，会话对象不支持序列化。长期保存的接收密钥仍采用原持久格式。该设计适合双方身份固定的连续通信，并带来明确的额外密钥存储成本。

\subsection{怎样计算连续消息的平均成本}
设一次会话准备耗时为$T_{\rm prep}$，会话内单条解密为$T_{\rm sess}$。处理$N$条消息的总成本与平均成本分别为
\[
 T_{\rm total}(N)=T_{\rm prep}+NT_{\rm sess},\qquad
 \bar T(N)=\frac{T_{\rm prep}}{N}+T_{\rm sess}.
\]
首条消息承担全部准备成本；连续消息逐渐摊薄这部分开销。第5章同时展示准备、后续解密和保留叶钥大小，以判断实际使用方式。
